% ============================================================================
%  RESULTS -- ONLINE SUPPLEMENT, SECTION 3: THE ESTIMATOR.
%  project `panel-npiv-discrete`
%
%  WHAT THIS FILE HOLDS.  Everything def:lfstep and def:crossfit carried that is
%  not the construction: the measurability verifications, the container
%  arguments, the re-declarations of conventions stated elsewhere, and the
%  discussion of what each clause rules out.  Copied verbatim from
%  results_s3_estimation.tex; nothing here was rewritten.  Only splice-point
%  punctuation and a lead-in clause were added where a moved fragment began or
%  ended mid-sentence, and each such fragment is headed by what it discusses.
%
%  IT IS READ AFTER results_s3_estimation.tex AND NOT INSTEAD OF IT.  The
%  paragraphs below cite the clauses they discuss, so the supplement inputs the
%  statements first.  Together the two files are the undivided Section 3 block.
%
%  NO \label IS DECLARED HERE.  Every labelled display of the two definitions --
%  eq:normconst, eq:cap, eq:lfsched, eq:prodmeas, eq:supnormconv, eq:foldrates --
%  stayed in the article, so no \eqref anywhere was orphaned and no equation
%  number moved.  This file contributes nothing to any counter.
%
%%  DO NOT \input THIS FILE ALONGSIDE results.tex.  results.tex survives
%%  undivided and remains the source of record.
% ============================================================================

\section*{Discussion of the estimator}

\noindent This section carries the material the article does not print about the two definitions that
construct the estimator: for the $\varphi$-step of \cref{def:lfstep} and the cross-fitting scheme of
\cref{def:crossfit}, the measurability verifications, the grounds on which each nuisance's container is
chosen, the re-declarations of conventions stated elsewhere, and the record of what each clause does and
does not rule out. The constructions themselves are printed above, verbatim and unchanged. Headings below
name what the text under them discusses.


% ------------------------------------------------------------------------
\subsection*{The $\varphi$-step (\cref{def:lfstep})}

% ------------------------------------------------------------------------
\paragraph{The measurability requirement on the preliminary, and where it is stated}

The opening of \cref{def:lfstep} declares the preliminary a measurable random element of $\mathcal V$
at $\sigma(\mathcal I_\ast)$, with $\mathcal I_\ast$ the fold fixed there: \emph{this is a requirement on the preliminary and
is stated \textbf{here}, where the preliminary is declared, rather than only inside the verification that
consumes it}; (P-a.1) of \cref{prp:phirate} carries the same requirement, and it is of the kind read off
a formula, which is why it costs nothing to impose.

% ------------------------------------------------------------------------
\paragraph{The sieve-membership clause, and why the $H^{s}$ bound is an envelope}

\emph{The sieve-membership clause is \textbf{imposed here}}, and (P-a.1) of \cref{prp:phirate}, this
opening's companion, carries it in terms. \emph{The clause is not idle}: every iterate after the zeroth
lies in $\mathcal S_n$
% ARCHAEOLOGY.  Through dispatch 20 this opening required measurability into V only
% and then imposed the H^s probability on the strength of a membership stated
% nowhere in it; the fourth bullet below establishes measurability of
% ||hat-varphi^{(j)} - varphi_dagger||_{H^s} expressly "at the configuration
% lem:pathcont states, where hat-varphi^{(0)} in S_n", a CONDITIONAL discharge of
% an UNCONDITIONAL requirement.  No result's hypothesis set changed.
by construction, so without it the preliminary is the one element of the path that need not, and the
$H^{s}$ display above is then a statement about a quantity no clause makes finite; \emph{envelopes, not $O_P$
statements, and \suppref{Remark}{paenvelope} is why}: the constants $\mathsf c_{14},\mathsf c_{15}$ below are
printed before the data are seen, and an $O_P$ hypothesis supplies no number for them to dominate.

% ------------------------------------------------------------------------
\paragraph{The letter $\omega$}

The letter is $\omega$ because \suppref{Lemma}{scaleinv}(i) shows the normalisation \textbf{is} a
Landweber relaxation parameter, which is what $\omega$ names in the regularisation literature.
% ARCHAEOLOGY.  The letter is not varsigma, which is a Sec.1.6 local bound variable
% and one subscript from the binding varsigma_0 of (P-a.2).  Glossary C67.

% ------------------------------------------------------------------------
\paragraph{What in the schedule is computable, and from what}

\emph{Every quantity in \eqref{eq:lfsched} is \textbf{computable} by the analyst, from four inputs and
two constants, and the two constants are not free}: $\hat\delta_n$ and $h_u$ are already inputs to the
procedure, $\kappa$ is a function of $\mu$, and $\hat\varepsilon_0$ is the declared rate of the
preliminary, those are the four; $\mathsf c_{14}$ and $\mathsf c_{15}$ are the two.
\emph{The scale \eqref{eq:normconst} and the cap \eqref{eq:cap} are a different matter, and the next
paragraph is where they are treated}: both read $\|1/\pi\|_\infty$, which no analyst knows.

% ------------------------------------------------------------------------
\paragraph{The reference density is unknown, what the floor buys, and the one thing it does not}

\Cref{def:spaces} declares $\pi$ an unknown feature of $P$ --- the pooled design density
$\tfrac12(f_{X_1}+f_{X_2})$ --- and \cref{def:lfstep} floors an estimate $\hat\pi$ of it at
$\varepsilon_\pi$. Three things follow, and the order matters, because the third is an open item and the
first two are what make it a narrow one.

\emph{(i) What the floor buys, and it is not only well-posedness --- nor, as this paragraph read
through 2026-08-26, only computability.} $\hat\pi\ge\varepsilon_\pi$ gives
$\|1/\hat\pi\|_\infty\le\varepsilon_\pi^{-1}$, so \suppeq{hatAadj} formed at $\hat\pi$ is defined at
every design point. \emph{The same floor also fixes the relaxation}, through
$\|h\|_{L^{2}(\mathcal X,dx)}\le\varepsilon_\pi^{-1/2}\|h\|_{L^{2}(\hat\pi)}$: that is the entry
$\varepsilon_\pi^{-1/2}$ in the maximum \eqref{eq:normconst} now carries, and by
\suppref{Proposition}{hatAprop}\textup{(i$'$)} it is what makes
$\|\omega_n^{-1}\hat{\mathcal A}_\varphi\|\le1$ true in the $L^{2}(\hat\pi)$ pairing the iteration
actually runs in --- not a substitute for an unknown $\|1/\pi\|_\infty^{1/2}$ but the constant
\textup{(F-f)} requires there. \emph{Through 2026-08-26 this paragraph offered exactly that number and
filed it under computability alone}, which is why the entry was absent from \eqref{eq:normconst} and
\suppref{Remark}{phirscope}\textup{(7)} enumerated four parts rather than five; the disposition is
\textup{(7e)}.
\emph{Computability then follows, and it now follows \textbf{on a range} rather than after a
substitution}: at $\varepsilon_\pi\le\|1/\pi\|^{-1}_\infty$ --- implied by the threshold clause
\textup{(ii)} names --- the maximum in \eqref{eq:normconst} \emph{is} $\varepsilon_\pi^{-1/2}$, leaving
$c_C$, $\mathsf c_2$ and $c_\Pi$, which C3(ii) and B1(iv.c) fix and the analyst therefore knows, and
$h_u,h_x$, which the analyst chooses; \eqref{eq:cap} then follows, being $\omega_n^{2}\hat\delta_n^{-2}$
at that $\omega_n$. \emph{Above the threshold the maximum reads $\|1/\pi\|_\infty^{1/2}$ and the two
displays are population displays and are \textbf{not} computable}, and that is worth saying beside the
paragraph above, which reports \eqref{eq:lfsched}
computable and is about a different set of quantities. \emph{Nor are those the only displays carrying
the norm}: \cref{rem:localise}\textup{(2)} names \suppref{Proposition}{hatAprop}\textup{(i)},
\eqref{eq:normconst}, \eqref{eq:cap}, \suppref{Lemma}{aerr} and \suppref{Lemma}{levels}\textup{(i)}, and
the Cauchy--Schwarz step of the measurability verification below is a further one. \emph{No claim of
exhaustiveness is made about either list.}

\emph{(ii) What the floor costs, and where the cost is charged.} The clipping is the metric projection
in $L^{2}(\mathcal X,dx)$ onto an order interval, and that set contains $\pi$ itself whenever
$\varepsilon_\pi\le\min\{\|1/\pi\|^{-1}_\infty,\|\pi\|^{-1}_\infty\}$; at such an $\varepsilon_\pi$ the
projection can only decrease $\|\hat\pi-\pi\|_{L^{2}(\mathcal X,dx)}$, which is
\cref{def:hatT}'s accounting for \eqref{eq:trunset} read one object over. C1(ii) makes that threshold
positive and names no value for it, so $\varepsilon_\pi$ joins $\mathsf c_{14}$, $\mathsf c_{15}$ and
$\tau_{\mathrm{DP}}$ as a constant the theory does not pin; above the threshold the clipping is a bias
this paper does not price. \emph{The population condition is untouched, and the direction of that
statement is the point}: C1(ii) bounds $\pi$ below on $\mathcal X$ and says nothing whatever about an
estimate of it, so a $\hat\pi$ that goes negative violates \textbf{no} condition the paper states --- it
is not a failure of C1(ii) and cannot be one. The convention in \cref{def:lfstep} is
\emph{additional} and \emph{sample-level}, and C1(ii) is neither weakened nor re-spelled to accommodate
it. \emph{The event is not hypothetical}: at the kernel order the estimator runs, $\hat\pi$ is a
higher-order density estimate, higher-order kernels are sign-changing, and the Monte Carlo of
Section~5 loses replications to exactly this failure when the operator is built on a small enough fold.

\emph{(iii) The open item.} What none of this establishes is that the substitution of $\hat\pi$ for
$\pi$ \emph{inside} the iteration leaves \eqref{eq:phirate} in place --- \emph{which is why
\cref{prp:phirate} states that display at $\hat\varphi^{\pi}$, the infeasible reading, from 2026-08-27,
and does not state it at $\hat\varphi$}.
\suppref{Remark}{phirscope}\textup{(7)} states that item in full, in five parts. \emph{Two are settled
outright}: the two norms are
equivalent with deterministic constants, so the norm \eqref{eq:phirate} is read in is not the
difficulty; and, as of 2026-08-27, the normalisation \eqref{eq:normconst} is admissible in \emph{both}
pairings, so \textup{(F-f)}'s $\|\mathcal L\|\le1$ holds whichever the iteration is read in --- that
last is \textup{(7e)}, and it was not among the parts this corpus enumerated before that date.
\emph{A third is settled \textbf{with a rider}, and the rider bears on the
verification below}: the feasible recursion is a genuine Landweber iteration on $L^{2}(\hat\pi)$, but
its sieve is a \textbf{random} subspace, because $\mathcal S_n$ is cut out
by $\pi$, whereas the four steps below are written at a fixed one; \emph{so those four steps verify the
measurability of $\hat\varphi^{\pi}$ and leave the measurability of $\hat\varphi$ pending}, which is
what \cref{def:crossfit}\textup{(c$'$.1)} now prints and what \suppref{Lemma}{epterm} consumes.
\emph{Two parts are simply open.} The level of the iterate is pinned by $\hat\pi$ and not by $\pi$, which
leaves a constant component of the error bounded by $\|\hat\pi-\pi\|_{L^{2}(\mathcal X,dx)}$ --- that
part is \emph{priced}, and what it needs is a rate condition on $\hat\pi$ of the form A4(iii) imposes on
$\hat\varrho$, which this paper does not impose. And
$\hat{\mathcal A}^\ast_{\hat\pi}=(\pi/\hat\pi)\hat{\mathcal A}^\ast_\varphi$ perturbs the operator
multiplicatively while (P-a.3) is a condition on the unperturbed one --- that part is not priced at all.
% ADDED 2026-08-26, PHASE 4d.  This paragraph exists because the deduction the
% theorist-critic recorded was SILENCE, not the absence of a theorem: the corpus
% divided by pi at eq:hatAadj, floored nothing, and said nothing about the gap.
% Saying "not established" in the corpus's own idiom is the deliverable; a lemma
% would be better and is not available without a uniform rate for pi-hat, which no
% assumption in this paper supplies.

% ------------------------------------------------------------------------
\paragraph{$\mathsf R$ against $\mathsf R_\ast$, and the two directions of a mis-set level}

The radius here is $\mathsf R_\ast$ and \textbf{not} $\mathsf R$, and the two must not be the same
symbol, this is the same species of defect as $\mathsf c_{14}$ versus $\mathsf c_{16}$ and it sits on
the same inequality. $\mathsf R$ is an \emph{$H^{s}$ ball radius}: it is what \suppref{Definition}{srcball}
measures, what A5(vi$'$) and (P-a.1) bound, and what (P-c), (P-g) and \suppeq{entropy} price, so it
carries units of Sobolev norm and enters \emph{approximation} statements. $\mathsf R_\ast$ is a
\emph{spectral pre-image norm} in $L^{2}(\pi)$: it is (P-a.3)'s bound on $\varsigma^{(k)}$, it enters
only through (F-f), and it is what $\mathsf c_{15}$ must dominate. They are not comparable in
either direction without the interpolation step \suppref{Remark}{pasrc}\textup{(1)} shows is unavailable.
The two directions of error are not symmetric. Chosen too large, the level is
conservative: the rule stops earlier, the rate in \eqref{eq:phirate} is unchanged in order and the
constant degrades. Chosen too small, the implemented level sits below the true data error
at some stage, (L3) fails, and with it (F-f): the discrepancy guarantee is void and the estimator is
whatever the cap \eqref{eq:cap} returns. \emph{This is the one place where the procedure carries a tuning constant the
% ARCHAEOLOGY.  Through round 11 one symbol carried both R and R_star, on the two
% sides of an inequality this definition has to establish; glossary C69.  The
% too-small failure mode is the one the fixed level tau_DP delta-hat_n ln n of
% round 9 produced, displaced from the schedule's SHAPE to its SCALE.
theory does not pin, and it is said here rather than left to be discovered}.

% ------------------------------------------------------------------------
\paragraph{The output is a measurable random element: why inspection of the formula will not settle it}

\emph{And the discharge
available for the other nuisances is not available here}: A4(iii) and \cref{def:crossfit}\textup{(a)} settle
$\hat\varrho$ ``by inspection of the formula'', which every kernel and every series construction passes;
inspection of the formula settles nothing for a Landweber--Fridman iterate stopped by a discrepancy rule,
\emph{because the stopping index of \textup{(c)} is random}.

% ------------------------------------------------------------------------
\paragraph{The verification, in four steps}

\emph{One index warning, because both readings occur in this paragraph}: the
superscript in $\hat\varphi^{(-k)}$ and $\mathcal I_{-k}$ just above is \cref{def:crossfit}'s \emph{fold}
index; every $k$ from here to the end of the verification is this definition's \emph{stage} index, and the
two are one glyph apart.
\begin{enumerate}
\item[--] \emph{The iterates.} Suppose $\hat\varphi^{(k-1)}$ is a $\sigma(\mathcal I_\ast)$-measurable random
element of $\mathcal V$. The finitely many numbers $\Delta_{h_x}\hat\varphi^{(k-1)}_i$, $i\in\mathcal I_\ast$, are
then random variables, being values of the \emph{bounded linear} functionals \suppeq{smdiff} at a measurable
random element, \emph{and the step from a random function to a random number is the
\textbf{Carath\'eodory} criterion and not the evaluation of a fixed functional, because the functional
\suppeq{smdiff} is itself \textbf{random}}, its kernel carrying $X_{i1},X_{i2}$ and $I_{i1},I_{i2}$, which lie in $\mathfrak I$ at every $n$ by C5$'$(ii) --- \emph{the per-$n$ form of that clause and not its limit}, since measurability is asserted at each $n$.
``Value of a bounded linear functional at a measurable random element'' gives measurability only for a
\emph{fixed} functional; what is available, and what is needed, is the superposition argument: the map
$F_i(\omega,\varphi):=\Delta_{h_x}\varphi_i(\omega)$ is $\sigma(\mathcal I_\ast)$-measurable in $\omega$ at each
fixed $\varphi$, because $X_{i1},X_{i2}$ are random variables and
$(x,x^\ast)\mapsto C_{h_x}\big(x-x^\ast,I(x^\ast)\big)=\Theta_{I(x^\ast)}\big((x-x^\ast)/h_x\big)$ is
\emph{jointly continuous on $\mathcal X\times\mathcal X$}, hence jointly Borel there.
\emph{The route is \textbf{not} continuity of $I\mapsto\Theta_I(\varsigma)$, and that map is
\textbf{not} continuous}: \suppref{Lemma}{bdryfam}'s proof of \textup{(d)} records that it jumps at
each $I$ whose endpoint crosses $\varsigma$, and \suppeq{ivalx} moves both endpoints of $I(x^\ast)$
affinely with $x^\ast$, so this is exactly the configuration that clause names. \emph{Nor is it
\textup{(f)} read as a statement about $\Theta$}: \textup{(f)} carries the \textbf{coefficient} map
$I\mapsto\mathsf G(I)^{-1}e_0$ and nothing more. \emph{What removes the jump is that the two arguments
move together, and the collapse is exact rather than approximate}: at $\varsigma=(x-x^\ast)/h_x$ and
$I=I(x^\ast)$ of \suppeq{ivalx}, $\varsigma\in I(x^\ast)$ holds if and only if
$\underline x\le x\le\bar x$, because $h_x>0$, so the indicator of \suppeq{gramfam} is
$\un\{x\in\mathcal X\}$ --- \emph{free of $x^\ast$}, and identically $1$ on the domain of integration
of \suppeq{smdiff}. On that domain the kernel is
$\big(\sum_{k}(\mathsf G(I(x^\ast))^{-1}e_0)_k\,\varsigma^{k}\big)e^{-\varsigma^{2}/2}$ with no
indicator left in it, and this is jointly continuous in $(x,x^\ast)$: $x^\ast\mapsto I(x^\ast)$ is
continuous into $\mathfrak I$ for the topology \suppeq{ifraktop}, both endpoints being affine in
$x^\ast$ and $\tanh$ continuous; $I\mapsto\mathsf G(I)^{-1}e_0$ is continuous by
\suppref{Lemma}{bdryfam}\textup{(f)}; and polynomial evaluation is jointly continuous in coefficient
vector and argument. So the $dx$-integral in \suppeq{smdiff} is measurable in $\omega$ by Tonelli
\emph{applied to positive and negative parts}, \emph{and the qualification is not decoration here}: the
integrand of \suppeq{smdiff} is a difference of two boundary-corrected kernels of order $p_K$, which
\eqref{eq:truncmom} makes \emph{sign-changing} by construction, against a $\varphi$ that integrates to zero
against $\pi$, so Tonelli in its nonnegative form does not apply to it as written. The split is legitimate
because the integrand is absolutely integrable at every sample point: $|C_{h_x}|\le\mathsf c_2$ by C3(ii)
and \textup{(P9)}, $\mathcal X$ is a bounded interval by A0, and
$\int_{\mathcal X}|\varphi|\,dx\le|\mathcal X|^{1/2}\|1/\pi\|^{1/2}_\infty\|\varphi\|_{L^{2}(\pi)}<\infty$
by C1(ii) and the Cauchy--Schwarz inequality, so each of the four one-signed integrals is finite and their
signed combination is well defined and measurable. \emph{This is the treatment the same paragraph gives the
$\Pi$-pairing below}; it is also
A2c\textup{(m.2)}'s and \cref{def:crossfit}\textup{(c$'$.1)}'s, on the other
side of the pairing. The integral is, moreover,
$\|\cdot\|_{L^{2}(\pi)}$-continuous in $\varphi$ at each fixed $\omega$, which is exactly the boundedness of
\suppeq{smdiff} that \suppref{Definition}{hatTexp} records and which holds at every sample point. Since $\mathcal V$
is a subspace of $L^{2}(\mathcal X,\pi)$ with $\mathcal X$ compact, it is separable, so the Carath\'eodory
criterion makes $F_i$ jointly $\sigma(\mathcal I_\ast)\otimes\mathfrak B(\mathcal V)$-measurable, and $F_i$ then
composes with the measurable map $\omega\mapsto\big(\omega,\hat\varphi^{(k-1)}(\omega)\big)$. \emph{The
device is this paper's own and is used at three further sites}: it is what \cref{def:crossfit}\textup{(c$'$.1)} names for
\eqref{eq:prodmeas}, what the opening of the proof of \cref{prp:frechet} runs for $G_\varphi$ and its
$u$-derivatives, and what C2$'$(i) runs for the envelopes $\bar f_k$;
% ARCHAEOLOGY.  Through dispatch 19 this link was stated as an evaluation, which is
% the one reading the randomness of eq:smdiff forecloses.
and $\hat T_{h_x}(\hat\varphi^{(k-1)},\theta)$, $\hat{\mathcal A}_{\hat\varphi^{(k-1)}}$ and
$\hat{\mathcal A}^\ast_{\hat\varphi^{(k-1)}}$ depend on $\hat\varphi^{(k-1)}$ through those numbers
\emph{alone}, by \suppeq{hatTx}, \suppeq{hatg} and \suppeq{hatAcon}, the adjoint being determined by the
operator and the two inner products of \cref{def:spaces}. Restricted to $\mathcal S_n$, which
is a \emph{fixed} finite-dimensional space, the iteration \textup{(b)} is therefore an affine recursion on
a Euclidean space of that dimension whose coefficients are random variables, started at the constant
$v_0=0$; so each
$v_m$, $m=0,1,\dots,\bar m_n$, is a $\sigma(\mathcal I_\ast)$-measurable random element of $\mathcal S_n$, being a
finite composition of measurable maps. That the coefficients are random variables is
\textbf{Tonelli} and not composition, and the distinction is the same one the previous link turns on.
\emph{The coordinates are those of a basis of $\mathcal S_n$ and \textbf{not} of $\{b_1,\dots,b_{L_x}\}$},
and two facts force the distinction:
$\mathcal S_n=\operatorname{span}(b_1,\dots,b_{L_x})\cap\mathcal V$, and the linear functional
$\zeta\mapsto\int_{\mathcal X}\zeta\,d\pi$ is \emph{not} identically zero on
$\operatorname{span}(b_1,\dots,b_{L_x})$: A8's $b_j$ are B-splines, hence \emph{nonnegative} and not
identically zero, and $\pi$ is bounded below away from zero by C1(ii), so
$\int b_j\,d\pi>0$ for every $j$, no $b_j$ lies in $\mathcal V$, $\{b_j\}$ spans a strictly larger space, and
$\dim\mathcal S_n=L_x-1$. \emph{Both consequences matter}: the coordinate map written in the $b_j$ is not a
coordinate map on $\mathcal S_n$, and $\hat{\mathcal A}_{\hat\varphi^{(k-1)}}b_j$ evaluates the operator off
the domain \cref{def:spaces} declares for it. \emph{Neither touches the conclusion}, and that is why the
repair is a rewording: fix once any basis of $\mathcal S_n$, every element of it lies in $\mathcal V$, so
the operator is applied only where it is defined, and every element is a finite linear combination of
$b_1,\dots,b_{L_x}$, and read the recursion in its coordinates. Its matrix is that of
$\hat{\mathcal A}^\ast_{\hat\varphi^{(k-1)}}\hat{\mathcal A}_{\hat\varphi^{(k-1)}}\vert_{\mathcal S_n}$,
whose entries are the $L^{2}(\Pi)$ pairings of the images under
$\hat{\mathcal A}_{\hat\varphi^{(k-1)}}$ of two basis elements, and its constant vector carries the
$L^{2}(\Pi)$ pairings of those images with $\hat T_{h_x}(\hat\varphi^{(k-1)},\theta)$
, each an \emph{integral over $(u,t,l)$ against $\Pi$ of a random integrand}, \suppeq{hatAadj} writing the
adjoint as exactly such an integral. \emph{A finite composition of measurable maps does not reach an
integral}; what does is joint measurability in $(\omega,u,t,l)$ followed by Tonelli's theorem. The joint
measurability is again Carath\'eodory, \emph{run at each contrast index separately}: by \suppeq{hatAcon}
and \suppeq{hatg}, at a fixed element $\zeta$ of $\mathcal V$ and a fixed
$l\in\{1,\dots,\mathsf L-1\}$ the maps
$(u,t)\mapsto\hat{\mathcal A}_{\hat\varphi^{(k-1)}}\zeta(u,t,l)$ and
$(u,t)\mapsto\hat T_{h_x}(\hat\varphi^{(k-1)},\theta)(u,t,l)$ are continuous at every sample point, $C_{h_u}$
and $\mathsf C_{h_u}$ being continuous in $u$ by \suppref{Lemma}{intker} and $t\mapsto w(t'\Phi(\Delta Z_i))$
continuous, while each is $\sigma(\mathcal I_\ast)$-measurable in $\omega$ at fixed $(u,t)$ by the paragraph
above, so both are $\sigma(\mathcal I_\ast)\otimes\mathfrak B(\IR\times\mathcal T)$-measurable at that $l$.
\emph{Measurability in $(\omega,u,t,l)$ then costs nothing, and the reason is that the third factor of
$\Pi$ carries a \textbf{finite} set}: $\{1,\dots,\mathsf L-1\}$ with its power set, so each map is the
finite disjoint union over $l$ of the pieces just shown measurable, and no further argument is needed.
Each factor lies
in $L^{2}(\Pi)$, by \suppref{Proposition}{hatAprop}(i) for the operator and, for $\hat T_{h_x}$, by
$\mathsf c_3c_\Pi<\infty$, a bound that reads the \emph{per-cell} $|\widehat{\widetilde W}_{i,l}|\le1$ of
\suppref{Remark}{hatTread}\textup{(2)} at each $l$ against a $\Pi_{\mathsf w}$ that B1(iv.d) makes a
\emph{probability} measure, and so does not move with $\mathsf L$, so the products are $\Pi$-integrable
by the Cauchy--Schwarz inequality, and Tonelli
applied to positive and negative parts then makes each pairing a $\sigma(\mathcal I_\ast)$-measurable function
of $\omega$. This is the inheritance
\cref{def:crossfit}\textup{(c$'$.1)} performs for $\omega\mapsto\|\hat\ell^{(-k)}-\ell_\star\|_{L^{2}(P)}$,
read on $\Pi$ instead of $P$. At $\mathsf L=2$ the third factor of $\Pi$ is the one-point set
$\{1\}$ and $\Pi_{\mathsf w}$ the point mass at it, so every display in this bullet is the one earlier
rounds printed.
\item[--] \emph{The stopping index.} $\tau_{\mathrm{DP}}$ and $\hat\delta^{\mathrm{eff}}_{n,k}$ are
deterministic by \eqref{eq:lfsched}, and $\|\hat{\mathcal A}_{\hat\varphi^{(k-1)}}v_m+\hat T_{h_x}(\hat\varphi^{(k-1)},\theta)\|_{L^{2}(\Pi)}$
is a random variable for each $m$, \emph{by the previous step \textbf{and Tonelli}, the two being
needed and neither sufficing alone}, since an $L^{2}(\Pi)$ norm is a $\Pi$-integral of a random function
and not a value of one: the previous step makes
$(\omega,u,t,l)\mapsto\big(\hat{\mathcal A}_{\hat\varphi^{(k-1)}}v_m+\hat T_{h_x}(\hat\varphi^{(k-1)},\theta)\big)(u,t,l)$
jointly measurable, $v_m$ being $\sigma(\mathcal I_\ast)$-measurable with values in the fixed
finite-dimensional $\mathcal S_n$, and Tonelli then makes $\omega\mapsto\int(\cdot)^{2}\,d\Pi$ measurable, the square root
being continuous. So, for each $m\le\bar m_n$, the event that \emph{that}
norm is at most $\tau_{\mathrm{DP}}\hat\delta^{\mathrm{eff}}_{n,k}$ is measurable, \emph{it carries no
letter}, every short candidate being claimed, and $\{\hat m^{(k)}=m\}$ is that event intersected with the
complements of the same events at the indices $m'<m$: a \emph{finite} intersection of measurable sets.
\emph{Two further features of \textup{(c)}, not the two the paragraph after this one names, which are
about the \emph{level}, are what this step consumes}: the cap
$\bar m_n$ of \eqref{eq:cap} is \emph{deterministic}, so the family is \textbf{finite} and no limit is taken;
and the emptiness convention makes $\hat m^{(k)}$ \textbf{total}, so no sample point is left without a value
and $\{\hat m^{(k)}=\bar m_n\}$ absorbs the residual event. A random cap or a partial index would break
the argument here, and \textbf{not} only here: the next bullet's finite sum of indicators needs the
same two features, and so does (L4) of \cref{prp:phirate}'s Step~2, which compares the stopping index
to $\bar m_n$ and then applies the emptiness convention on the complement. \emph{No claim of
exhaustiveness is made, and that is deliberate}: the draft this sentence replaces read ``and nowhere
else'' and was short by its own next bullet.
\item[--] \emph{The output.} $v_{\hat m^{(k)}}=\sum_{m\le\bar m_n}\un\{\hat m^{(k)}=m\}\,v_m$ is a
\emph{finite} sum of indicators of measurable events against measurable random elements, hence a
$\sigma(\mathcal I_\ast)$-measurable random element of $\mathcal S_n$; and $\hat\varphi^{(k)}$ of \textup{(d)} is
one of $\mathcal V$, $\mathcal V$ being a linear space and $\mathcal S_n\subset\mathcal V$
(\cref{def:spaces}). Induction on $k\le\bar k$, from a $\sigma(\mathcal I_\ast)$-measurable
$\hat\varphi^{(0)}$, \emph{which is a requirement on the preliminary, and \textbf{is} of the kind read off a
formula; it is \textbf{stated} in this definition's opening and in (P-a.1) of \cref{prp:phirate}, and is
cited rather than imposed here}, gives the claim for $\hat\varphi=\hat\varphi^{(\bar k)}$.
\item[--] The $H^{s}$ norm of the same object, which the paragraph owes because \suppref{Lemma}{pathcont} and
A5(vi$'$) are $H^{s}$ statements. Both take probabilities of
$\max_{j\le\bar k}\|\hat\varphi^{(j)}-\vdagger\|_{H^{s}}\le\mathsf R$, so
$\omega\mapsto\|\hat\varphi^{(j)}-\vdagger\|_{H^{s}}$ must be a random variable as well, and measurability
into the $L^{2}(\pi)$ container does not by itself give it. \emph{It follows from the three steps above at the
configuration \suppref{Lemma}{pathcont} states}, where $\hat\varphi^{(0)}\in\mathcal S_n$ and hence the whole path
lies in $\mathcal S_n$: since $\mathcal S_n\subset H^{s}$ by \eqref{eq:sieveinhs}, $\|\cdot\|_{H^{s}}$ is a
norm on a \emph{fixed} finite-dimensional space, hence a continuous, and so measurable, function of the
coordinates, which the steps above make random variables. \emph{The membership is A8's business and it is now
\textbf{stated} there rather than flagged here}: nothing in this definition places the B-spline basis in
$H^{s}$, and \eqref{eq:sieveinhs} is A8's basis clause read as the containment it delivers.
\emph{\suppeq{jackson} is \textbf{not} the right pointer, and the error would be one of direction
rather than of precision}: \suppeq{jackson} states how well splines \emph{approximate}
$H^{s}$ functions, and what is needed is that splines \emph{lie in} $H^{s}$. \emph{Approximation power does
% ARCHAEOLOGY.  Through dispatch 19 this bullet pointed at eq:jackson.
not give membership, and the two conditions are genuinely different inequalities on the same degree}, a
basis of degree $d$ approximates $H^{s}$ at the rate \suppeq{jackson} states as soon as $s\le d+1$, while
$\mathcal S_n\subset H^{s}$ requires $d>s-\tfrac12$; on the window $s-1\le d\le s-\tfrac12$ the first holds
and the second fails, and the witness $s=1.7$, $d=1$ sits in it. A8's basis clause is the
right home, as this bullet's own first sentence already said, and it is stated there in \emph{degree}, the
one reading of ``order'' on which the spline conventions agree.
\end{enumerate}
This is a discharged verification and \textbf{not} a fifth entry in \suppref{Remark}{obligations}'s inventory,
and the distinction is that remark's own. That inventory carries obligations whose \emph{discharge site is a
later step}, something a stated result's truth depends on and no present result establishes. Here the
discharge is performed in the same definition that constructs the object, so nothing is owed forward.
\emph{That is not the same as having no consumer, and the draft this sentence replaces said it
was}: \cref{def:crossfit}\textup{(c$'$.1)} consumes exactly this measurability when it fixes the
jointly measurable representative of $\hat\varphi^{(-k)}$, which is what supplies the composite at
\suppref{Lemma}{prelimtheta} Steps~4 and~4$'$ --- the second precedent named just below is therefore
also a consumer. The precedents are A4(iii)'s joint-measurability half and
\cref{def:crossfit}\textup{(c$'$.1)}'s Tonelli inheritance, neither of which is inventoried either, for the
same reason.

% ------------------------------------------------------------------------
\paragraph{The two load-bearing features of the stopping rule, and the status of Bakushinskii's veto}

Two features of (c) are load-bearing and one of them is a correction. First, the level is
stage-dependent, and it must be. The data error of the stage-$k$ equation (a) is not $\hat\delta_n$
but $\hat\delta_n\ln n$ \emph{plus} the linearisation remainder $h_u^{-2}\varepsilon_{k-1}^{2}$, where
$\varepsilon_{k-1}=\|\vdagger-\hat\varphi^{(k-1)}\|_{L^{2}(\pi)}$, see \suppeq{stagedecomp}. A
fixed level $\tau_{\mathrm{DP}}\hat\delta_n\ln n$ sits
\emph{below} the true data error in exactly the regime the recursion must pass through, namely
$h_u^{-2}\varepsilon_{k-1}^{2}\ge\hat\delta_n\ln n$; there the classical discrepancy guarantee does not
apply and the rule either over-fits the noise or fails to stop before $\bar m_n$.
Second, the level is deterministic and supplied from outside, and the reason has two parts of
very different status. The first is this paper's own and is sharp; the second is classical and cited,
and is weaker than it looks.
% ARCHAEOLOGY.  Rounds through 9 implemented the fixed level tau_DP delta-hat_n ln n,
% a mismatch between the estimator and its analysis; it is repaired at the estimator
% rather than in the proof.  A round-9 draft over-read Bakushinskii's veto twice.

(2a) The realised residual cannot supply the level, and the reason is what ill-posedness is. The
natural data-driven repair is to replace $h_u^{-2}\hat\varepsilon_{k-1}^{2}$ by
$h_u^{-2}\hat r_{k-1}^{2}$ with $\hat r_{k-1}$ the residual realised at stage $k-1$. \emph{It cannot
work.} The quantity the level must dominate is $h_u^{-2}\varepsilon_{k-1}^{2}$ with
$\varepsilon_{k-1}=\|\vdagger-\hat\varphi^{(k-1)}\|_{L^{2}(\pi)}$, and a residual bounds $\varepsilon_{k-1}$
from below only through a lower bound on $\|\hat{\mathcal A}_{\varphi_\circ}v\|$ in terms of $\|v\|$.
That $\|\hat{\mathcal A}v_\dagger\|\ll\|v_\dagger\|$ is available is not an accident of this
construction; it is the definition of ill-posedness. So $\hat r_{k-1}$ can be arbitrarily smaller than
$\varepsilon_{k-1}$, and a threshold built from it sits below the data error for exactly the reason the
fixed threshold did. This observation is the paper's own and it is the sharp part of the
argument.

(2b) The general statement is Bakushinskii's veto, it is cited, and it is weaker than it looks in
a statistical setting. That \emph{no} rule depending on the data alone, and not on the noise level
, can make a regularisation method convergent for a genuinely ill-posed problem is
\citet{bakushinskii1984}; see \citet[\S3.3]{engl2000}. Two scope statements, because the veto is
routinely over-read. It is a \emph{worst-case, deterministic}
statement: convergence is required uniformly over the admissible data, and the counterexample it turns
on is a worst case. \emph{In the statistical setting this paper inhabits it does not rule out
data-driven selection}: Lepskii-type balancing over the regularisation index is purely data-driven and
attains the optimal rate up to a logarithmic factor \citep{lepskii1990}, and Section~5 runs such a
stop alongside the schedule. \emph{So the schedule is preferred here, not forced here}, and the reasons are three and
local. The number of stages $\bar k$ is finite and $\bar k_a$ is deterministic, so there is no
adaptation problem across stages to solve. A deterministic level is what makes the rule
\cref{def:lfstep}(c) implements and the rule \cref{prp:phirate} Step~2 analyses coincide \emph{by
construction}. And a Lepskii stop over $m$ would
need a variance proxy at each candidate index that the operator-noise model of this paper does not
supply without a further bound. A balancing stop is therefore a live alternative rather than a
foreclosed one, and were it to work here the requirement that (P-a.2) deliver a \emph{rate} would
soften materially.


% ------------------------------------------------------------------------
\subsection*{The $\theta$-step (\cref{def:crossfit})}

% ------------------------------------------------------------------------
\paragraph{Convention \textup{(P8)}: the transfer rule and its two companions}

\textup{(P8)}'s \emph{transfer rule} and
its two companions, the $O_P$ companion and the $o_P$ companion, are \suppref{Lemma}{pconv}(i)--(iii),
stated at a \emph{generic} $\sigma$-field, written $\mathcal G$ there and instantiated here, and they apply
with no extension.
% ARCHAEOLOGY.  Through dispatch 16 this clause read that the generic sigma-field
% was "unnamed there", which was true of the comment block those statements then
% lived in and is not true of the compiled lemma.

% ------------------------------------------------------------------------
\paragraph{The conditioning $\sigma$-field is written out and given no symbol}

The licence is the same
either way, since the convention is a reading of a proof and not a hypothesis; the sentence is not, and this
definition is where the re-declaration lives. \emph{That $\sigma$-field is written out at every site and is
given no symbol}: the
three natural candidates are all claimed, $\mathcal G$ is \suppref{Lemma}{density}'s closed span,
$\mathcal E$ is the best-approximation error of \suppeq{jackson}, and $\mathcal F$ is the Fourier
transform of \eqref{eq:usobolev}, and \suppref{Lemma}{balance} and \textup{(P8)} already write
$\sigma(\mathcal I_a)$ out, so the convention is inherited rather than invented.

% ------------------------------------------------------------------------
\paragraph{How each conditional bound of the $\theta$-step is produced, and which companion converts it}

Each conditional bound below is produced by conditional Markov or conditional Chebyshev, hence holds at
\emph{every} threshold with a $\sigma(\mathcal I_{-k})$-measurable numerator, and is converted by the
\emph{$O_P$ companion} when the numerator is unconditionally $O_P(1)$ after division by the stated rate
(\suppref{Lemma}{prelimtheta} Step~4, \suppref{Lemma}{sandwich} Step~2) and by the \emph{$o_P$ companion} when it is
unconditionally $o_P(1)$ (\suppref{Lemma}{epterm} Step~6, \suppref{Lemma}{sandwich} Step~8). \emph{The transfer rule
needs the conditional probability to vanish at a \textbf{single} moving threshold, which neither
inequality delivers}; that is the observation \textup{(P8)} already recorded for the $o_P$ case, now
applied to the $O_P$ case. What supplies the tightness of the numerators is A4(iii), A5(ii), (iv) and
hypothesis (R-a) of \cref{thm:rootn}.

% ------------------------------------------------------------------------
\paragraph{What clause \textup{(c)}'s three properties buy}

The three properties clause \textup{(c)} requires of $\hat\ell^{(-k)}$ are the minimum
\cref{def:weight}(b) needs to be well posed, they are what make $\int_0^u\hat\ell^{(-k)}(v,z)dv$ exist
and make $\partial_u\hat L^{(-k)}=\hat\ell^{(-k)}$ hold at every point, \emph{and they are stated here
rather than inherited}, so that this definition is well posed on its own terms.

% ------------------------------------------------------------------------
\paragraph{Clause \textup{(c$'$)} re-declares nothing: what \suppref{Lemma}{balance} already states}

\emph{This clause re-declares nothing}:
\suppref{Lemma}{balance}'s sample-splitting convention states
$\sigma(\mathcal I_a)$-measurability of $\hat\ell$ as a \emph{random element} of
$L^{2}(P_{\Delta U,\Delta Z})$ together with a fixed \emph{jointly measurable representative}, which is
strictly weaker than \eqref{eq:prodmeas} below and is all that site needs, it evaluating $\hat\ell$ at an
\emph{undistorted} $(\Delta U,\Delta Z)$ draw.
% ARCHAEOLOGY.  Through dispatch 15 lem:balance's convention stated MEMBERSHIP and
% a measurability whose target was unspecified, and this sentence in turn described
% it as "a membership statement and not a measurability declaration".  Both were
% repaired in dispatch 16 at lem:balance.

% ------------------------------------------------------------------------
\paragraph{Why the carve-out for $\hat\varphi^{(-k)}$ is forced}

For $\hat\varphi^{(-k)}$ the property of clause \textup{(c$'$.1)} is verified rather than required; \emph{and the exception is forced, not stylistic}, since
one cannot require a property of an object one has already built, \suppref{Remark}{obligations}'s test, read on an
estimator.
% ARCHAEOLOGY.  Through dispatch 18 this sentence quantified over all five without
% the carve-out, so the clause simultaneously pinned hat-varphi^{(-k)} to a
% construction and demanded of it a property that construction was never shown to
% have.

% ------------------------------------------------------------------------
\paragraph{Clause \textup{(c$'$.1)}: what supplies $\hat\varphi^{(-k)}$'s measurability, and the choice of representative}

Clause \textup{(c$'$.1)} records that the ground which settles the other two, A4(iii)'s ``by inspection
of the formula'', is unavailable for $\hat\varphi^{(-k)}$: no formula is inspected when an index is
chosen by a
discrepancy rule. \emph{What supplies it instead is the measurability paragraph of \cref{def:lfstep}}, which
builds the property from the deterministic cap \eqref{eq:cap}, the emptiness convention of
\cref{def:lfstep}\textup{(c)}, and the finiteness of the resulting family: the stopping index is a hitting
index of finitely many measurable events, the iterates are finite compositions of measurable maps, and the
stopped iterate is a finite indicator-weighted sum of them. \emph{The superscripts are not written out here},
because \cref{def:lfstep}'s index is a \emph{stage} index and this clause's is a \emph{fold} index, and the
two are one glyph apart. \emph{The other two remain requirements, as printed}, and $\hat b^{(-k)}$'s is read
off its formula on the $\hat\varrho$ model.
\emph{The sample $\sigma$-algebra is written out and given no symbol}, because
calligraphic $\mathcal F$ is A2c's Fourier transform and $\mathcal F_{XZW}$ never appears
without its subscript; the same convention applies to $\sigma(\mathcal I_{-k})$
itself.
% ARCHAEOLOGY.  Glossary Sec.1.6, C15, and the C78 write-it-out precedent.
For these three the choice of representative is immaterial, and that is the test a container has to
pass, but the test is passed on \textbf{two different grounds}. For $\hat\varrho^{(-k)}$ and $\hat b^{(-k)}$ the evaluation is at
$\Delta Z_i$, whose law \emph{is} the one the container is built on, so identity of laws does it. \emph{For
$\hat\varphi^{(-k)}$ it is not}: its container is $\mathcal V\subset L^{2}(\mathcal X,\pi)$, whose null
sets are the $\pi$-null ones, while its evaluations are at $X_{i1}$ and $X_{i2}$, whose laws are $f_{X_1}$
and $f_{X_2}$. \emph{What transfers the null sets there is \textbf{domination}, and it has a tag}: C1(ii)
gives $\bar c_\pi=\max_{j\in\{1,2\}}\|f_{X_j}/\pi\|_\infty<\infty$, so every $\pi$-null set is
$P_{X_j}$-null for $j\in\{1,2\}$; \emph{``the law its container is built on'' is the wrong ground for the
third object}. Either way any two representatives
% ARCHAEOLOGY.  The domination ground was uncited through dispatch 15.
agree at all of the finitely many evaluation points outside a $P$-null set, the union of the finitely many
exceptional sets being $P$-null.

% ------------------------------------------------------------------------
\paragraph{Clause \textup{(c$'$.1)}: the inheritance is Tonelli on positive and negative parts}

\emph{The inheritance is the
\textbf{measurability half} of Fubini--Tonelli and it is taken on positive and negative parts,
entrywise}: \eqref{eq:weightdef}'s integrand is $\IR^{p\times(\mathsf L-1)}$-valued and
% ARCHAEOLOGY.  A dispatch-15 draft elided this to "by Tonelli's theorem applied to
% int_0^u hat-ell^{(-k)}(v,z) dv".
of arbitrary sign, and $\int_0^u=-\int_u^0$ when $u<0$. For each entry $(j,l)$ with $j\le p$ and
$l\le\mathsf L-1$, and each sign part, the
two maps
\[
(\omega,u,z)\;\longmapsto\;\int_\IR\un\{0\le v\le u\}\,\big(\hat\ell^{(-k)}_{jl}\big)^{\pm}(\omega)(v,z)\,dv,
\qquad
(\omega,u,z)\;\longmapsto\;\int_\IR\un\{u\le v\le0\}\,\big(\hat\ell^{(-k)}_{jl}\big)^{\pm}(\omega)(v,z)\,dv
\]
are product measurable by Tonelli's theorem, their integrands being non-negative and
$\sigma(\mathcal I_{-k})\otimes\mathfrak B(\IR\times\IR\times\IR^{p})$-measurable by \eqref{eq:prodmeas};
all four are \emph{finite} at every point because clause \textup{(c)} makes $\hat\ell^{(-k)}$ bounded at
each realisation, so their signed combination is well defined and product measurable. \emph{This is the
device A2c\textup{(m.2)} already owns}, on the other side of the
pairing, where the same positive-and-negative-part split carries Tonelli across a signed
$u$-integral.

% ------------------------------------------------------------------------
\paragraph{Clause \textup{(c$'$.1)}: why the derivative inherits product measurability under (R-a)}

\emph{Under (R-a) the derivative inherits it too, and \textup{(c$'$.2)} needs that for its second
envelope}: $\partial_u\hat\ell^{(-k)}$ exists at every point under (R-a) and is the pointwise limit
of the difference quotients of $\hat\ell^{(-k)}$ in $u$ taken along the reciprocals of the positive
integers, \emph{the increment is written out and carries no letter}, on the same ground as
\eqref{eq:supnormconv}'s countable dense set; each such quotient is product measurable by
\eqref{eq:prodmeas}, and a pointwise limit of countably many product-measurable functions is product
measurable. \emph{A representative of an $L^{2}$ class supplies nothing of the kind}, having no derivative
in general.
No equivalence class is formed and no representative is chosen, so there is nothing
for a choice of representative to disturb, and clause \textup{(c)}'s continuity in $u$ is untouched.

% ------------------------------------------------------------------------
\paragraph{Clause \textup{(c$'$.1)}: why the $L^{2}$ reading fails for the two weight objects}

\emph{Why the $L^{2}$ reading fails here and not three lines above}: these two are evaluated at a
% ARCHAEOLOGY.  A draft of this clause declared the two weight objects random
% elements of L^2(P_{Delta U, Delta Z}) and asserted that "the continuity clause of
% (c) is preserved by the choice", which is false for an arbitrary representative,
% point evaluation not being continuous on L^2.
\emph{shifted} first argument, at $(\Delta U-\lambda\Delta h_k,\Delta Z)$ in \suppref{Lemma}{biasexp} Steps~1
and~4--5, at $(\Delta U_i(\theta,\hat\varphi^{(-k)}),\Delta Z_i)$ in \eqref{eq:Jhat} and
\suppref{Lemma}{sandwich} Steps~2 and~9, whose law \textup{(c$'$.2)} below argues, and needs, \emph{not} to be
dominated by $P_{\Delta U,\Delta Z}$; on an undominated set two representatives of one class disagree, so
$\hat J$ of \eqref{eq:Jhat}, $\Psi^{(1)}_k$ of \suppref{Lemma}{biasexp}, and hence $\hat\theta$ of
\eqref{eq:onestep} and $\hat\Sigma$ of \eqref{eq:Sighat}, would depend on which was chosen and
\emph{the estimator would not be a function of the object the clause declares}.

% ------------------------------------------------------------------------
\paragraph{Clause \textup{(c$'$.1)}: the derivation of $\hat L^{(-k)}$'s membership}

The membership of $\hat L^{(-k)}$ in $L^{2}(P_{\Delta U,\Delta Z})$ is derived and not required: by \eqref{eq:weightdef},
$\|\hat L^{(-k)}(u,z)\|\le\|\hat b^{(-k)}(z)\|+\|\hat\ell^{(-k)}\|_\infty|u|$, so
$E\|\hat L^{(-k)}(\Delta U,\Delta Z)\|^{2}\le
2E\|\hat b^{(-k)}(\Delta Z)\|^{2}+2\|\hat\ell^{(-k)}\|^{2}_\infty E[\Delta U^{2}]<\infty$ almost surely
under \cref{def:weight}(b)'s admissibility, (R-a) and A1b(iv). Both statements are about the
function and neither is a container declaration.

% ------------------------------------------------------------------------
\paragraph{Clause \textup{(c$'$.1)}: the seven \textbf{Part-I} sites of the freezing lemma, and the direction the proofs use}

\emph{Equation \eqref{eq:prodmeas} is what \textup{(P8)}'s freezing lemma consumes for the
$\hat\ell^{(-k)}$ \emph{ingredient} of its composite, at six of its seven \textbf{Part-I} sites}, that lemma being
stated at a composite of the form $(\omega,e)\mapsto f\big(\hat\ell^{(-k)}(\omega),e\big)$. \emph{The
seventh is \suppref{Lemma}{prelimtheta} Steps~4 and~4$'$, where $\hat\ell^{(-k)}$ does not appear at all and the
frozen object is $h_k$}: what supplies the composite there is the jointly measurable representative fixed
above for $\hat\varphi^{(-k)}$, so both halves of this clause are load-bearing and neither covers the
other. \emph{Seven is the \textbf{Part-I} count and not the count of the lemma's sites, which is
nine}, and the qualifier is load-bearing because the other two are supplied by something else:
\suppref{Lemma}{opnoise} Step~4 sits in Part~H, is stated before \cref{def:crossfit} exists and does
not adopt it, and takes its composite from A4(iii); \suppref{Lemma}{balance} Steps~0 and~2 sit in
Part~F and take theirs from that lemma's own sample-splitting convention, which is deliberately
weaker than \eqref{eq:prodmeas}. The catalogue is \suppref{Remark}{pfreezesites}, which carries all
nine; this paragraph points into it and maintains no count of its own. The composite is \textbf{not} product measurable ``precisely when'' \eqref{eq:prodmeas}
holds, in either direction. Not sufficient: the composite is product
measurable only once $f$ is itself jointly Borel in its two arguments, and at the sites where the
evaluation point is built from $\hat\varphi^{(-k)}$ only together with that object's own clause. Not
necessary: $f\equiv0$ is product measurable whatever $\hat\ell^{(-k)}$ is. \emph{What is true is the
direction the proofs use}, and it is the one this clause is for: \eqref{eq:prodmeas} is what this
% ARCHAEOLOGY.  A dispatch-15 draft wrote the "precisely when".
definition \emph{supplies}, and it is what the per-realisation ``Borel'' of clause \textup{(c)} does not. \emph{The file owns the device by which \eqref{eq:prodmeas} is verified, and it
is recorded rather than imposed}: \suppref{Lemma}{biasexp} Step~1 obtains joint measurability in
$(\lambda,\omega)$ from ``continuous in $\lambda$, measurable in $\omega$'' by the Carath\'eodory
criterion, and the same criterion gives \eqref{eq:prodmeas} whenever, for each fixed $u$, the map
$(\omega,z)\mapsto\hat\ell^{(-k)}(\omega)(u,z)$ is
$\sigma(\mathcal I_{-k})\otimes\mathfrak B(\IR^{p})$-measurable, since \textup{(c)} already makes
$u\mapsto\hat\ell^{(-k)}(u,z)$ continuous at every $(\omega,z)$. \emph{That is a sufficient condition and
not the requirement}; the requirement is \eqref{eq:prodmeas}, which is what the proofs use.

% ------------------------------------------------------------------------
\paragraph{Clause \textup{(c$'$.2)}: measurability of the outer essential supremum}

The \textbf{outer} essential supremum is measurable as well, and this does \textbf{not} follow from
joint measurability without an argument, it is not a countable supremum. The argument is the
tool of the previous
paragraph, once more. At each fixed real level, the map
% ARCHAEOLOGY.  A dispatch-15 draft of this clause carried the outer essential
% supremum on the word "likewise".
$\omega\mapsto P_{\Delta Z}\big\{z:\sup_u\|\hat\ell^{(-k)}(\omega)(u,z)\|>\text{that level}\big\}$ is
$\sigma(\mathcal I_{-k})$-measurable by Tonelli's theorem applied to the non-negative jointly measurable
indicator; the levels at which that probability vanishes are exactly the levels \emph{at or above} the
essential supremum, since it vanishes at the essential supremum itself, the exceptional set there being
the union of the countably many exceptional sets one reciprocal-of-an-integer above it, and at no
smaller level, by the definition of the essential supremum as an infimum; and therefore, for every real
level, the event that the essential supremum does not exceed it is the intersection, over the
\emph{countably many rational levels strictly above it}, of the events that the probability vanishes, 
a countable intersection of measurable sets, the \emph{strict} inequality being what makes those
exceptional sets exhaust the one at the level itself.

% ------------------------------------------------------------------------
\paragraph{Clause \textup{(c$'$.2)}: why the envelope's measurability is load-bearing}

\emph{The levels are written out and carry no letter}, on the same ground as the countable
dense set above. \emph{This is load-bearing and not hygiene}: (R-a) prices both envelopes as
$O_P(1)$, which is a statement about the distribution of these two maps, and \textup{(P8)}'s $O_P$
companion requires its \emph{bounding numerator} to be measurable with respect to the conditioning
$\sigma$-field, and at \suppref{Lemma}{sandwich} Step~2 that numerator is
$\|\hat\ell^{(-k)}\|^{2}_\infty E\|\Delta Z\|^{2}$, whose only random factor is the envelope. \emph{No
symbol of the companion's is imported into compiled text}; the numerator is named by what it is.

% ------------------------------------------------------------------------
\paragraph{Clause \textup{(c$'$.2)}: a genuine supremum over $z$ would not be measurable}

\emph{A genuine supremum over $z$, by contrast, would not be a measurable function of $\omega$}, since no
clause makes $z\mapsto\hat\ell^{(-k)}(u,z)$ continuous.

% ------------------------------------------------------------------------
\paragraph{Clause \textup{(c$'$.2)}: where the exceptional set falls, and why $\ell_\star$ needs no convention}

\emph{The exceptional set is therefore $P$-null where the envelope is used}:
for each $i$, $\sup_u\|\hat\ell^{(-k)}(u,\Delta Z_i)\|\le\|\hat\ell^{(-k)}\|_\infty$ almost surely, and
the union over the finitely many $i\le n$ of the exceptional sets is still $P$-null, so the bound holds
simultaneously at all $n$ observations almost surely. \emph{For the deterministic $\ell_\star$ no
convention is needed and none is imposed}: A$2'$(i)'s closing clause defines $\ell_\star$ pointwise and
\suppref{Corollary}{elldiff} bounds it at \emph{every} $(u,z)$, so $\|\ell_\star\|_\infty$ is a genuine supremum and
the two readings agree.

% ------------------------------------------------------------------------
\paragraph{Clause \textup{(d)}: what the second display would say without the sieve premise}

The second display of \eqref{eq:foldrates} is finite at every realisation only under
\eqref{eq:sieveinhs}: without the first the display is an $O_P(1)$ statement about a quantity equal to
$+\infty$ surely, which is a condition on the \emph{empty} event and not a rate.

% ------------------------------------------------------------------------
\paragraph{Clause \textup{(d)}: what the fold reading of A5(vi$'$) costs}

This costs nothing and is nonetheless not without consequence. It costs nothing because $\mathsf F$
is fixed, so the fold reading is the intersection of $\mathsf F$ events each of probability tending to one;
it is not without consequence because the radius $\mathsf R$ must then serve all $\mathsf F$ fold
complements \emph{at one value}, $\mathsf R$ being a constant of \suppref{Definition}{srcball} and not a fold-indexed
object, and $\mathsf R$ also appears in (P-c), (P-g), \suppeq{entropy} and elsewhere, so a
fold-indexed radius is not available. \emph{Those three are named as instances and this paragraph
maintains no list}: \suppref{Definition}{srcball} carries the enumeration and this points at it. The
schedule constants carry $\mathsf R_\ast$, which that definition rules a \textbf{different} radius,
not comparable to $\mathsf R$ in either direction. (T-5) already reads
A5(vi$'$) at each fold \emph{operationally}, through ``\eqref{eq:phirate} holds on each
$\mathcal I_{-k}$''; what was missing is that this enumeration, which declares itself the place where
the reading lives, did not list it. \emph{That operational reading survives the 2026-08-27 scoping of
\eqref{eq:phirate} and is worth saying so}: what (T-5) delivers at each fold is the display at
$\hat\varphi^{\pi}$, and A5(vi$'$) is a hypothesis \emph{consumed} inside \cref{prp:phirate} at
\suppref{Lemma}{pathcont}(c) rather than a conclusion drawn from it, so the fold reading it needs is of
the hypothesis and is untouched by which reading of \cref{def:lfstep}\textup{(b)} the conclusion names.
\emph{What does \textbf{not} survive is the parallel reading for A5(ii)}, and that clause is carried at
\textup{(T-4)}.


% ------------------------------------------------------------------------
\paragraph{The parts of \textup{(P8)} do not have the same standing}

% SPLICE ARTIFACT REPAIRED 2026-08-09.  The body printed its own heading twice,
% the second time after a bare em-dash with no verb following:
%   "...do not have the same standing.--- and the parts of (P8) do not have the
%    same standing"
% Deduplicated.  BUT NOTE WHAT REMAINS: the paragraph now says only what its
% heading says.  The explanation the heading promises -- WHICH parts, and how
% their standing differs -- is NOT ON THE PAGE and was presumably lost in the
% 2026-08-02 splice that produced the duplication.  This is a CONTENT gap in
% printed supplement text, not a typographical one, and it is recorded as such
% rather than papered over with a plausible-sounding sentence.  Owner: whoever
% next opens this file with the (P8) history to hand.
Convention \textup{(P8)} is re-declared at \cref{def:crossfit}, and the parts of
\textup{(P8)} do not have the same standing.

% ------------------------------------------------------------------------
\paragraph{How clause \textup{(c$'$)}'s two sub-clauses stand to clause \textup{(c)}}

Two clauses follow; neither is available from clause \textup{(c)} alone, and
\textup{(c$'$.1)} is a requirement on the estimators and not a consequence of anything above it, 
\emph{with \textbf{one} exception, marked at the clause itself, and it is the object clause \textup{(b)} has
already fully constructed}.

% ------------------------------------------------------------------------
\paragraph{Clause \textup{(a)}'s truncation is a pointwise operation on the very function it names}

For $\hat\varrho^{(-k)}$ no representative is chosen at all; \cref{def:hatT}'s truncation is then a pointwise operation on the very
% SPLICE ARTIFACT REPAIRED 2026-08-09: the sentence ended ", and." -- truncated
% mid-clause.  The dangling conjunction is removed and the sentence closed at the
% clause that was complete.  What followed "and" is not recoverable from the page;
% if the paragraph owed a second consequence, it is still owed.
function clause \textup{(a)} names.

% ------------------------------------------------------------------------
\paragraph{Why the $L^{2}$ container is nonetheless retained for $\hat\varrho^{(-k)}$}

The $L^{2}$ container is retained
for $\hat\varrho^{(-k)}$ only because A4(iii)'s \emph{rate} is an $L^{2}$ statement; it is not what any
pointwise bound rests on.
